% jm.ever@schaersvoorde.nl %
\documentclass[dutch,12 pt,english]{article}
\usepackage[latin9]{inputenc}
\usepackage{geometry}                                                                                                                      
\geometry{verbose,a4paper,tmargin=0.5cm,bmargin=0.5cm,lmargin=1.2cm,rmargin=2.1cm,headheight=0.5cm,headsep=0.1cm,footskip=0.8cm}               
\pagestyle{empty} 
\usepackage{amsmath}
\usepackage{amssymb}
\usepackage[T1]{fontenc}
\usepackage{babel}
\usepackage{graphics}
\usepackage{graphicx}
\usepackage{mathptmx} % nodig als Font voor ACROBAT
\usepackage{helvet}
\usepackage{amsmath} % denk aan in mathmode: \xrightarrow{heele lange text boven rechtsepijl}
\usepackage{multicol}
\usepackage{multirow}
\usepackage{marvosym} %denk aan euro: \EUR \EURcr \EURtm
\usepackage{color}
\usepackage{eepic} %denk aan \circle{diameter}
\usepackage[scanall]{psfrag}
%\usepackage[figure]{hypcap}
\usepackage[colorlinks=true, urlcolor=blue,linkcolor=blue, pdftitle={GRM en Kansrekening},pdfauthor={J.M. Evers}, bookmarksnumbered=true,pdfnewwindow=true, pdfstartview=XYZ,pdfpagelayout=OneColumn,bookmarksopen]{hyperref}

\newcommand{\noun}[1]{\textsc{#1}}
\AtBeginDocument{ \def\labelitemii{\ding{54}} }
\makeatletter
\makeatother

\begin{document}
    
    \begin{flushleft}
    
    \begin{center}
	\noun{\LARGE Bord van Galton } \
	
	
	\noun{Empirische en theoretische kans }
	
    \end{center}
    \footnote{\includegraphics[  width=1cm, height=1cm]{./plaatjes/linux.jpg}                                 
    \emph{werd gemaakt onder} $LinuX$ \emph{met} \LaTeX 
    
    J.M. Evers 9/11/2009
    }




\indent


Hier volgt een voorbeeld opgave.

\indent

\indent

\begin{tabular}{|l|}
\hline
In een bord van Galton laat je tenminste 65 knikkers in 6 bakjes vallen.\tabularnewline 

De kans voor een bal om naar rechts te gaan is 0.3.\tabularnewline

Bepaal aan de hand van je resultaten:\tabularnewline

- de empirische kans\tabularnewline

- de theoretische kans\tabularnewline

\hline
\end{tabular}

\indent

\indent


\noun{Empirische kans }

Een empirische kans is een proefondervindelijke kans. \\
Deze bereken je dus nadat je een proef hebt gedaan. \\
In dit voorbeeld moet je dan 65 knikkers laten vallen in het bord van Galton.



\begin{multicols}{2}
Voor de proef

\includegraphics[scale=0.30]{./plaatjes/galton1.jpg}

Na de proef

\includegraphics[scale=0.30]{./plaatjes/galton2.jpg}

\end{multicols}

De empirische kans bereken je dan alsvolgd:

\indent

\begin{tabular}{|l c c c c c c |}
    \hline
    Nummer bakje  & 1 & 2 & 3 & 4 & 5 & 6\tabularnewline
    \hline
    Aantal  & 15 & 17 & 22 & 10 & 1 & 0\tabularnewline
    \hline
     & & & & & &  \tabularnewline
    Empirische kans & $\frac{15}{65}$&$\frac{17}{65}$ &$\frac{22}{65}$ & $\frac{10}{65}$&$\frac{1}{65}$ & 0   \tabularnewline
    \hline
\end{tabular}

\newpage

\noun{Theoretische kans }

\indent


Bij het berekenen van een theoretische kans heb je geen proef nodig. 

Je moet goed kijken hoe vaak er een 'keuze' wordt gemaakt tussen 'links vallen' en 'rechts vallen'.\\
In dit voorbeeld wordt er 5 keer een 'keuze' gemaakt.\\
Om in bakje 6 te komen valt de knikker 5 keer naar rechts.


\indent

\begin{multicols}{2}
Stel X staat voor het aantal keren rechts.
X is binomiaal verdeeld met n=5 en p=0.3

$P(X=0)=0.7^5 \approx 0.1681$\\
$P(X=1)=\left( \begin{array}{c}5 \\ 1 \end{array} \right) 0.3 \cdot 0.7^4 \approx 0.3602$\\
$P(X=2)=\left( \begin{array}{c}5 \\ 2 \end{array} \right) 0.3^2 \cdot 0.7^3 \approx 0.3087$\\
$P(X=3)=\left( \begin{array}{c}5 \\ 3 \end{array} \right) 0.3^3 \cdot 0.7^2 \approx 0.1323$\\
$P(X=4)=\left( \begin{array}{c}5 \\ 4 \end{array} \right) 0.3^4 \cdot 0.7 \approx 0.0.0284$\\
$P(X=5)=0.3^5 \approx 0.0024$

\includegraphics[scale=0.30]{./plaatjes/galton5.jpg}

\end{multicols}

\indent

\begin{tabular}{|l c c c c c c |}
    \hline
    Nummer bakje  & 1 & 2 & 3 & 4 & 5 & 6\tabularnewline
    \hline
    Aantal rechts  & 0 & 1 & 2 & 3 & 4 & 5 \tabularnewline
    \hline 
    Theoretische kans & 0.1681 & 0.3602 & 0.3087 & 0.1323 & 0.0284 & 0.0024 \tabularnewline
    \hline 
\end{tabular}

\indent

Je kunt de theoretische kansverdeling ook mbv je GRM uitrekenen.\\
Met TI \\
y1=binompdf(5,0.3,X) 

De kansen kun je vervolgens aflezen in de tabel.

\indent

Dit is het bijbehorende wims-plaatje.

\includegraphics[scale=0.70]{./plaatjes/galton4.jpg}






 \end{flushleft}
\end{document}

% TIPS TIPS TIPS TIPS TIPS TIPS TIPS TIPS %   
%  \href{http://wims.math.leidenuniv.nl/wims/wims.cgi?lang=nl&+module=H4/stat/stat-2.nl/cmd=new/subject=9/level=1/total_exos=1/rounding=-1/usage=2/taal=nl}{kans berekenen} 
%  \'{o} ó
%  `{o} ò
%  \^{o} ô
%  "{o} ö
% bold font : \textbf{ deze tekst is Bold }
% italic font: \emph{ deze tekst is schuin... }
%
% afdwingen: \begin{samepage} \end{samepage} let op: wel goed " nesten " bijvoorbeeld binnen een enumerate blijven
%
% nieuwe regel: 2 x backslash  \%
%
% tabel basis:
% de colommen: {ccc} of {lll} of {c|c|c} of {|c|c|c|}
% de breedte voor multiline in een cel:  {c|c|p{breedte cm}}
% de breedte voor multiline in een cel:  {p{3cm} | p{4cm} | p{5cm}}
% gebruik bij een vaste celbreedte: \newline om een nieuwe regel in de cel af te dwingen

% \begin{tabular}{c|c|p{3cm}} 
%	kop 1 & kop2  & kop3   \tabularnewline                                                                    
% 	\hline                                                                                                             
% 	text 1 & text 2  & text 3a \newline text3b \newline text3c   \tabularnewline                                                                    
%	text 1 & text 2  & text 3   \tabularnewline                                                                    
% 	text 1 & text 2  & text 3   \tabularnewline                                                                    
% 	text 1 & text 2  & text 3   \tabularnewline                                                                    
% \end{tabular} 

% plaatjes: \includegraphics[scale=0.3]{ ./VWO5.wiA.p2.2008/1228922811.plaatjes/plaatje.jpg }
% plaatjes: \includegraphics[width=10cm,height=4cm]{ ./VWO5.wiA.p2.2008/1228922811.plaatjes/plaatje.jpg } 
% plaatjes: \includegraphics[width=10cm,keepaspectratio]{ ./VWO5.wiA.p2.2008/1228922811.plaatjes/plaatje.jpg }

% gebruik in mathmode \displaystyle als je grotere breuken/wortels wilt
%
% TIPS TIPS TIPS TIPS TIPS TIPS TIPS TIPS %    

% JOUW PLAATJES:

% plaatje 1 : copy&paste  \includegraphics[scale=0.3]{./1228922811.plaatjes/helling.jpg}

%	 \includegraphics[scale=0.2]{./plaatjes/normaal.jpg}
%	 
%	 \includegraphics[scale=0.2]{./plaatjes/normaal1.jpg}
	 
%	 \includegraphics[scale=0.2]{./plaatjes/normaal2.jpg}
	 
%	 \includegraphics[scale=0.2]{./plaatjes/normaal3.jpg}
	 
%	 \includegraphics[scale=0.2]{./plaatjes/normaal4.jpg}
	
% plaatje 1 : copy&paste  \includegraphics[scale=0.3]{./plaatjes/galton1.jpg}
% plaatje 2 : copy&paste  \includegraphics[scale=0.3]{./plaatjes/galton2.jpg}
% plaatje 3 : copy&paste  \includegraphics[scale=0.3]{./plaatjes/galton3.jpg}
% plaatje 4 : copy&paste  \includegraphics[scale=0.3]{./plaatjes/galton4.jpg}
% plaatje 1 : copy&paste  \includegraphics[scale=0.3]{./plaatjes/galton5.jpg}
